\documentclass[10pt]{article}
\usepackage[letterpaper,margin=0.78in]{geometry}
\usepackage{amsmath,amssymb,amsthm}
\usepackage{microtype}
\usepackage{xcolor}
\usepackage[colorlinks=true,linkcolor=blue!55!black,urlcolor=blue!55!black,citecolor=blue!55!black]{hyperref}
\hypersetup{
  pdftitle={The necessity half of Deddens's intertwining conjecture},
  pdfauthor={fa banach 001, agent lane 00, model GPT5.6}
}

\newtheorem{theorem}{Theorem}
\newtheorem{lemma}[theorem]{Lemma}
\newcommand{\D}{\mathbb D}
\newcommand{\Hp}{H^2}
\newcommand{\spt}{\sigma_{\mathrm p}}

\title{Removing the kernel zeros in Deddens's\\
analytic-Toeplitz intertwining conjecture}
\author{Autonomous mathematical discovery run \texttt{fa\_banach\_001}\\
\texttt{agent\_lane\_00} --- model GPT5.6}
\date{13 August 2026}

\begin{document}
\maketitle

\begin{abstract}
Bourdon and Shapiro record as open the necessity half of Deddens's
conjecture: if two analytic Toeplitz operators satisfy
$XT_\varphi=T_\psi X$ for some nonzero bounded $X$, must
$\overline{\psi(\D)}$ lie in the point spectrum of $T_\varphi^*$?  We prove
that it must.  The usual reproducing-kernel argument gives an analytic
eigenvector field except where $X^*$ annihilates a kernel.  At each such
point, factor out the finite-order zero of this Hilbert-space-valued analytic
field.  The nonzero leading coefficient remains an eigenvector with the
missing eigenvalue.
\end{abstract}

\paragraph{Status.}
This is a candidate new full proof of the necessity statement declared open
in the introduction of \cite{BourdonShapiro}.  A bounded search found no later
resolution; no priority claim is made.

\section{The source question}

For $\varphi,\psi\in H^\infty(\D)$, let $T_\varphi$ and $T_\psi$ denote
multiplication by the symbols on the Hardy space $\Hp$.  Deddens conjectured
that
\begin{equation}\label{eq:deddens}
 \overline{\psi(\D)}\subseteq\spt(T_\varphi^*)
\end{equation}
is necessary and sufficient for the existence of a nonzero bounded operator
$X$ such that $XT_\varphi=T_\psi X$.  Here the bar means pointwise complex
conjugation, not closure.  Cowen disproved sufficiency; Bourdon and Shapiro
state that necessity remains open.  Their kernel argument proves the desired
point-spectrum assertion away from a discrete zero set and only a
full-spectrum assertion at the remaining points.

\begin{theorem}\label{thm:main}
Let $\varphi,\psi\in H^\infty(\D)$.  If a nonzero bounded operator
$X:\Hp\to\Hp$ satisfies
\[
 XT_\varphi=T_\psi X,
\]
then \eqref{eq:deddens} holds.  Thus the necessity half of Deddens's
conjecture is true, without additional assumptions on either symbol.
\end{theorem}

\section{Idea of the proof}

The adjoint of an analytic Toeplitz operator has a canonical eigenvector at
every point: the reproducing kernel $K_a$ has eigenvalue
$\overline{\psi(a)}$ for $T_\psi^*$.  Applying $X^*$ produces a
$T_\varphi^*$-eigenvector unless $X^*K_a=0$.  These exceptional points are
the zeros of a nonzero Hilbert-space-valued analytic function.

At an exceptional point, do not discard the eigenvector field.  Divide it by
the largest common power of the local coordinate.  The intertwining identity
survives off the zero, hence by continuity also at the zero; the value of the
divided field there is nonzero and is precisely the missing eigenvector.

\section{An analytic zero-removal lemma}

\begin{lemma}\label{lem:zero-removal}
Let $A$ be a bounded operator on a complex Banach space $E$, let
$F:\D\to E$ be analytic and not identically zero, and let
$a:\D\to\mathbb C$ be analytic.  If
\begin{equation}\label{eq:eigenfield}
 AF(\zeta)=a(\zeta)F(\zeta)\qquad(\zeta\in\D),
\end{equation}
then $a(\D)\subseteq\spt(A)$.
\end{lemma}

\begin{proof}
Fix $\zeta_0\in\D$.  If $F(\zeta_0)\ne0$, equation
\eqref{eq:eigenfield} gives the assertion.  Otherwise, the Taylor expansion
of the nonzero Banach-valued analytic function $F$ has a first nonzero
coefficient at some finite order $m\ge1$.  Thus locally
\[
 F(\zeta)=(\zeta-\zeta_0)^mG(\zeta),
\]
where $G$ is analytic and $G(\zeta_0)\ne0$.  Cancelling the scalar factor in
\eqref{eq:eigenfield} for $\zeta\ne\zeta_0$ gives
$AG(\zeta)=a(\zeta)G(\zeta)$.  Continuity at $\zeta_0$ yields
\[
 AG(\zeta_0)=a(\zeta_0)G(\zeta_0),
 \qquad G(\zeta_0)\ne0.
\]
Hence $a(\zeta_0)$ is an eigenvalue.  Since $\zeta_0$ was arbitrary, the
lemma follows.
\end{proof}

\section{Resolution of the conjectured necessity}

\begin{proof}[Proof of Theorem~\ref{thm:main}]
For $\zeta\in\D$, define the Hardy-space vector
\[
 k_\zeta(w)=\frac{1}{1-\zeta w}\qquad(w\in\D).
\]
The map $\zeta\mapsto k_\zeta$ is $\Hp$-valued analytic; in coefficient
coordinates it is $(1,\zeta,\zeta^2,\ldots)$.  Put
\[
 F(\zeta)=X^*k_\zeta,
 \qquad
 \psi^\sharp(\zeta)=\overline{\psi(\overline\zeta)}.
\]
Both functions are analytic.  Moreover, $F$ is not identically zero.  Indeed,
$k_\zeta=K_{\overline\zeta}$, and the span of all Hardy reproducing kernels
is dense; if $X^*$ killed every $k_\zeta$, then $X^*=0$ and hence $X=0$.

Taking adjoints in the intertwining identity gives
$T_\varphi^*X^*=X^*T_\psi^*$.  Since
\[
 T_\psi^*k_\zeta=\psi^\sharp(\zeta)k_\zeta,
\]
we obtain the analytic eigenfield identity
\[
 T_\varphi^*F(\zeta)=\psi^\sharp(\zeta)F(\zeta)
 \qquad(\zeta\in\D).
\]
Lemma~\ref{lem:zero-removal} therefore implies
$\psi^\sharp(\D)\subseteq\spt(T_\varphi^*)$.  Finally,
for each $a\in\D$ take $\zeta=\overline a$ to get
\[
 \psi^\sharp(\overline a)=\overline{\psi(a)}
 \in\spt(T_\varphi^*).
\]
This is exactly \eqref{eq:deddens}.
\end{proof}

\section*{Verification notes}

The conjugate parameter $\zeta=\overline a$ is used solely to make the
kernel field analytic.  The zero-removal step is valid for arbitrary
Banach-valued analytic functions: nonidentity to zero guarantees a first
nonzero Taylor coefficient at every zero.  The proof establishes point
spectrum, not merely spectrum, because the divided field is nonzero at every
formerly exceptional point.  The source question page and a separate audit
accompany this packet.

\begin{thebibliography}{9}
\bibitem{BourdonShapiro}
P.~S. Bourdon and J.~H. Shapiro,
\emph{Intertwining relations and extended eigenvalues for analytic Toeplitz
operators}, Illinois J. Math. 52 (2008), 1007--1030;
arXiv:0801.1972.
\end{thebibliography}

\end{document}
